% Les sept composés carbonylés saturés de formule brute C5H10O (corrigé du défi D08)
{\renewcommand{\arraystretch}{1.1}\setlength{\tabcolsep}{3pt}%
\setchemfig{atom sep=1.5em}%
\newcommand{\nomI}[1]{{\otLight\fontsize{7.5}{9}\selectfont\color{otMuted}#1}}%
\begin{tabular}{*{3}{>{\centering\arraybackslash}m{46mm}}}
\multicolumn{3}{l}{\otMedium\fontsize{8}{10}\selectfont\color{otNavy}Aldéhydes} \\
\chemfig{CH_3-CH_2-CH_2-CH_2-CHO} & \chemfig{CH_3-CH_2-CH(-[6]CH_3)-CHO} & \chemfig{CH_3-CH(-[6]CH_3)-CH_2-CHO} \\
\nomI{pentanal} & \nomI{2-méthylbutanal} & \textbf{\nomI{3-méthylbutanal}} \\[1mm]
\chemfig{CH_3-C(-[2]CH_3)(-[6]CH_3)-CHO} & & \\
\nomI{2,2-diméthylpropanal} & & \\[2mm]
\multicolumn{3}{l}{\otMedium\fontsize{8}{10}\selectfont\color{otNavy}Cétones} \\
\chemfig{CH_3-CO-CH_2-CH_2-CH_3} & \chemfig{CH_3-CH_2-CO-CH_2-CH_3} & \chemfig{CH_3-CO-CH(-[6]CH_3)-CH_3} \\
\nomI{pentan-2-one} & \nomI{pentan-3-one} & \nomI{3-méthylbutan-2-one} \\
\end{tabular}}
\par\smallskip\legende{Formules semi-développées des sept composés carbonylés saturés de formule $\ce{C5H10O}$.}
